\documentclass[11pt,a4paper]{article}
\usepackage{amsmath,amssymb,amsthm}
\usepackage{geometry}
\usepackage{hyperref}
\usepackage{booktabs}
\usepackage{parskip}

\geometry{margin=2.5cm}

\newcommand{\E}{\mathbb{E}}
\newcommand{\N}{\mathcal{N}}
\newcommand{\Lcal}{\mathcal{L}}
\newcommand{\q}{\mathbf{q}}
\newcommand{\x}{\mathbf{x}}
\newcommand{\R}{\mathbb{R}}

\title{%
  Position-Resolved Likelihood Inference\\
  for Atom Gradiometer Image Data%
}
\author{}
\date{}

\begin{document}
\maketitle
\tableofcontents
\newpage

%%============================================================
\section{Setup and Notation}
%%============================================================

\subsection{Gradiometer Architecture}

We consider an atom gradiometer composed of two atom interferometers (AIs)
operated simultaneously, one with atoms sourced at height $z = 0$\,m and one at
$z = 100$\,m, indexed hereafter by $z \in \mathcal{Z} = \{0,\,100\}$.
Each AI independently measures the accumulated phase along its atom trajectories.
Common-mode laser-phase noise cancels in the differential signal, while a
spatially-dependent differential phase (e.g.\ from a gravitational wave or
dark-matter field) is retained.

A run consists of $N_\mathrm{shots}$ shots indexed $i = 1,\ldots,N_\mathrm{shots}$,
acquired at times $t_i$.  Each shot produces two 2D fluorescence images per AI:
one for ground-state atoms ($s = g$, bright port) and one for excited-state atoms
($s = e$, dark port).  Images share a common pixel grid with bin centres
$\x_b = (x_b, y_b)$, bin area $\Delta A = \Delta x\,\Delta y$, and
$N_b$ bins total.

\subsection{Parameters}

\paragraph{Science signal.}
The target quantity is the differential phase signal injected into AI $z=100$:
\begin{equation}
  \delta\phi_i(\beta)
  = A_s \sin(2\pi f\,t_i) + A_c \cos(2\pi f\,t_i),
  \qquad \beta = (A_s,\,A_c),
  \label{eq:signal}
\end{equation}
where the signal frequency $f$ is assumed known a priori, leaving two free
parameters.  The signal amplitude and phase are
$A = \sqrt{A_s^2 + A_c^2}$ and $\varphi = \arctan2(A_c, A_s)$.

\paragraph{Common laser phase.}
The common laser phase $\phi_i \sim \mathrm{Uniform}(0, 2\pi)$ is shared by
both AIs within each shot, drawn independently across shots, and unobserved.
It is a latent variable; how we handle it in the likelihood is discussed in
Section~\ref{sec:nuisances}.

\paragraph{Effective AI phase.}
The total phase entering the transition probability for AI $z$ in shot $i$ is
\begin{equation}
  \Phi_{iz}(\phi_i,\beta)
  = \phi_i + \psi_z\!\bigl(\delta\phi_i(\beta)\bigr),
  \qquad
  \psi_0 = 0,\quad \psi_{100} = \delta\phi_i(\beta).
  \label{eq:eff_phase}
\end{equation}

\paragraph{Cloud nuisance parameters.}
The initial atomic cloud for AI $z$ in shot $i$ is characterised by a
Gaussian phase-space distribution with parameters
\begin{equation}
  \eta_{iz}
  =
  \bigl(
    \mu_{x_0},\,\mu_{y_0},\,\mu_{v_x},\,\mu_{v_y},\,
    \sigma_x,\,\sigma_y,\,\sigma_{v_x},\,\sigma_{v_y}
  \bigr)_{iz},
\end{equation}
the transverse centre-of-mass (COM) positions, COM velocities, and Gaussian
widths.  These vary shot-to-shot and across AIs.  Like $\phi_i$, they are
unobserved and must be handled carefully in the likelihood; see
Section~\ref{sec:nuisances}.

%%============================================================
\section{Semiclassical Phase-Space Map (PSMAP)}
%%============================================================

\subsection{PSMAP Surrogate}

The semiclassical simulation (ais++) integrates the interferometer sequence for a
regular 4D grid of initial transverse phase-space coordinates
$\q_0 = (x_0, y_0, v_{x0}, v_{y0})$ at fixed source height and longitudinal
velocity.  For each grid point and each output port $p$ it stores:
(i) wavepacket path amplitudes $a_{0,p}(\q_0)$ and $a_{1,p}(\q_0)$;
(ii) accumulated phase difference $\Delta\phi_p(\q_0)$;
(iii) interference flag $\iota_p \in \{0,1\}$; and
(iv) output state label $s_p \in \{g,e\}$.
At runtime, values are obtained by quadrilinear interpolation on the
$21^4$ grid (\texttt{PSMAPSurrogate}).
A separate PSMAP is computed for each AI.

\subsection{Transition Probabilities}

The probability that an atom with initial condition $\q_0$ is detected at port
$p$ given total AI phase $\Phi$ is
\begin{equation}
  P_p(\q_0;\,\Phi)
  = a_{0,p}^2 + a_{1,p}^2
    + \iota_p \cdot 2\,a_{0,p}\,a_{1,p}\,
      \cos\!\bigl(\Delta\phi_p(\q_0) + \Phi\bigr).
  \label{eq:port_prob}
\end{equation}
The per-state probabilities are obtained by summing over ports with matching state
label:
\begin{equation}
  p_s(\q_0;\,\Phi)
  = \sum_{p:\,s_p = s} P_p(\q_0;\,\Phi),
  \qquad s \in \{g, e\}.
  \label{eq:state_prob}
\end{equation}
Because the simulation retains only the two dominant interferometry paths,
\begin{equation}
  p_g(\q_0;\,\Phi) + p_e(\q_0;\,\Phi) \leq 1;
  \label{eq:pruning}
\end{equation}
atoms not accounted for are lost to pruned paths.  In particular,
$p_e \neq 1 - p_g$ in general; both are computed independently.

\subsection{Ballistic Trajectory Map}

Under free-flight propagation over detection time $T_\mathrm{det}$,
the final transverse position is
\begin{equation}
  \x^f(\q_0)
  =
  \begin{pmatrix}
    x_0 + T_\mathrm{det}\,v_{x0} \\
    y_0 + T_\mathrm{det}\,v_{y0}
  \end{pmatrix}.
  \label{eq:ballistic}
\end{equation}
Both output states follow the same trajectory (no state-dependent force before
imaging).  With default parameters ($T_\mathrm{det} = 3.8$\,s,
$\sigma_x = 100\,\mu$m, $\sigma_{v_x} = 309\,\mu$m/s), the cloud expands to
$\sigma_{x^f} \approx 1.2$\,mm at detection.

%%============================================================
\section{Per-Pixel Statistics via QMC}
\label{sec:qmc}
%%============================================================

\subsection{ACS Decomposition}

Expanding the cosine in~\eqref{eq:port_prob}, the phase dependence of $p_s$ is
entirely captured by three $\Phi$-independent scalars per pixel $b$, AI $z$, state $s$:
\begin{align}
  A_{zsb}(\eta_{iz})
  &= \E_{\q_0 \sim \text{cloud},\, \x^f\in b}
     \!\bigl[a_{0}^2(\q_0) + a_{1}^2(\q_0)\bigr],
  \label{eq:Ab} \\[4pt]
  C_{zsb}^c(\eta_{iz})
  &= \E_{\q_0 \sim \text{cloud},\, \x^f\in b}
     \!\bigl[2\,\iota\,a_{0}(\q_0)\,a_{1}(\q_0)\cos\Delta\phi(\q_0)\bigr],
  \label{eq:Cc} \\[4pt]
  C_{zsb}^s(\eta_{iz})
  &= \E_{\q_0 \sim \text{cloud},\, \x^f\in b}
     \!\bigl[-2\,\iota\,a_{0}(\q_0)\,a_{1}(\q_0)\sin\Delta\phi(\q_0)\bigr],
  \label{eq:Cs}
\end{align}
where the expectation is over initial conditions $\q_0$ drawn from the Gaussian
cloud $\N(\eta_{iz})$ whose final positions $\x^f(\q_0)$ fall within pixel $b$.
The expected pixel count is then
\begin{equation}
  \boxed{
  \lambda_{izsb}(\Phi)
  \;=\;
  L_{iz}\;
  \bigl[A_{zsb} + C_{zsb}^c\cos\Phi + C_{zsb}^s\sin\Phi\bigr],
  }
  \label{eq:lambda_analytic}
\end{equation}
where $\Phi = \Phi_{iz}(\phi_i,\beta)$ from~\eqref{eq:eff_phase} and $L_{iz}$
is an atom-number scale factor (Section~\ref{sec:atom_number}).

\subsection{QMC Evaluation (\texttt{SurrogatePixelACS})}

The conditional expectations~\eqref{eq:Ab}--\eqref{eq:Cs} are evaluated by
quasi-Monte Carlo (QMC) quadrature with $n_\mathrm{quad}$ Sobol points.

\begin{enumerate}
  \item Draw $n_\mathrm{quad}$ quasi-random standard-normal deviates
        $\mathbf{z}_j \in \R^4$, $j=1,\ldots,n_\mathrm{quad}$, using a
        scrambled Sobol sequence.  These are fixed at construction time and
        reused across all $\theta$ evaluations.

  \item Transform to cloud Gaussian: $\q_{0,j} = \mu(\eta) + \mathrm{diag}(\sigma(\eta))\,\mathbf{z}_j$.

  \item Compute detection positions $\x^f_j$ via~\eqref{eq:ballistic} and
        assign each sample to pixel $b_j$ using pixel edge arrays.
        Samples outside the image boundary are discarded.

  \item Evaluate the PSMAP at all $\q_{0,j}$ simultaneously via vectorised
        quadrilinear interpolation on the \texttt{PSMAPSurrogate} grid.
        This yields amplitudes $(a_{0,jp}, a_{1,jp})$ and phase residuals
        $\Delta\phi_{jp}$ for all ports $p$, in one GPU kernel with no Python loops.

  \item Read off the ACS components directly from the PSMAP outputs:
        \begin{equation}
          A_{jp} = a_{0,jp}^2 + a_{1,jp}^2, \quad
          C^c_{jp} = \iota_p\,2\,a_{0,jp}\,a_{1,jp}\cos\Delta\phi_{jp}, \quad
          C^s_{jp} = -\iota_p\,2\,a_{0,jp}\,a_{1,jp}\sin\Delta\phi_{jp}.
        \end{equation}
        Sum over ports with matching state label to obtain
        $(A_j, C^c_j, C^s_j)$ per sample per state.

  \item Accumulate per-pixel ACS by equal-weight bincount:
        \begin{equation}
          A_{zsb} = \frac{1}{n_\mathrm{in}}
                    \sum_{j:\,b_j = b} A_{j},
        \end{equation}
        and analogously for $C^c_{zsb}$, $C^s_{zsb}$, where
        $n_\mathrm{in}$ is the number of samples inside the detector.
        Both ground and excited state statistics are accumulated simultaneously.
\end{enumerate}

The default $n_\mathrm{quad} = 20{,}000$ Sobol points (rounded up to the next
power of 2) gives ACS tensors with sub-percent statistical noise in occupied pixels,
at a cost of one vectorised PSMAP evaluation per call.

%%============================================================
\section{Generative Model and Likelihood}
%%============================================================

\subsection{Poisson Observation Model}

Conditional on $\phi_i$ and $\eta_{iz}$, pixel counts are independent Poisson:
\begin{equation}
  n_{izsb} \;\sim\; \mathrm{Poisson}\!\bigl(\lambda_{izsb}(\Phi_{iz})\bigr),
  \qquad s \in \{g,e\},\; b = 1,\ldots,N_b.
  \label{eq:poisson}
\end{equation}
The per-shot conditional log-likelihood (for both AIs jointly) is
\begin{equation}
  \log\Lcal_i(\beta,\phi_i,\{\eta_{iz}\})
  = \sum_{z,s,b}
    \Bigl[
      n_{izsb} \log \lambda_{izsb}(\Phi_{iz})
      - \lambda_{izsb}(\Phi_{iz})
    \Bigr],
  \label{eq:cond_ll}
\end{equation}
where $\log(n!)$ constants are dropped as parameter-free.

\subsection{Atom Number Scale Factor}
\label{sec:atom_number}

Rather than tracking the launched atom count $\Lambda_{iz}$ explicitly, we
absorb it into a per-shot, per-AI scale factor $L_{iz}$ defined by the
requirement that predicted total counts match observed total counts at the
profiled cloud parameters $\hat\eta_{iz}$:
\begin{equation}
  L_{iz}
  = \frac{\displaystyle\sum_b \bigl(n_{izgb} + n_{izeb}\bigr)}
         {\displaystyle\sum_b \bigl(A_{zsb}(\hat\eta_{iz})
                              + A_{zsb}^{(e)}(\hat\eta_{iz})\bigr)},
  \label{eq:Lfix}
\end{equation}
where the denominator sums the $\phi$-independent terms over all pixels for
both states.  This normalization is exact whenever the interference terms
$C^c$ and $C^s$ integrate to zero over all pixels (complementarity), which
holds to high accuracy in practice.  $L_{iz}$ is computed once at
$\hat\eta_{iz}$ and treated as fixed.

%%============================================================
\section{Handling Unobserved Parameters: Profiling and Marginalisation}
\label{sec:nuisances}
%%============================================================

The full likelihood~\eqref{eq:cond_ll} involves two kinds of unobserved parameters
per shot: the common laser phase $\phi_i$, and the cloud parameters $\eta_{iz}$.
Both must be eliminated before optimising over $\beta$.

\subsection{Marginalise $\phi_i$}

$\phi_i \sim \mathrm{Uniform}(0, 2\pi)$ is completely unknown.  Profiling
(jointly optimising $\phi_i$ with $\beta$) would allow the per-shot phase to
absorb the signal entirely, biasing $\hat\beta$ toward zero.  We therefore
marginalise numerically using a uniform phase grid
$\phi_k = 2\pi k / K$, $k = 0,\ldots,K-1$:
\begin{equation}
  \log\Lcal_i^\mathrm{marg}(\beta,\{\hat\eta_{iz}\})
  = \operatorname{logsumexp}_{k=0}^{K-1}
    \!\bigl[\log\Lcal_i(\beta,\phi_k,\{\hat\eta_{iz}\})\bigr]
    - \log K.
  \label{eq:lse}
\end{equation}
This approximates the exact integral over $\phi_i$ to arbitrary accuracy as
$K \to \infty$.  The required grid size scales as
\begin{equation}
  K \;\gtrsim\; 2\sqrt{\bar N},
  \label{eq:ntheta}
\end{equation}
where $\bar N$ is the mean detected atom count per shot per AI.  In practice
$K = 512$ is used.

\subsection{Profile $\eta_{iz}$}

Each $\eta_{iz}$ has 8 components; for a $32\times32$ image, there are 2048
Poisson measurements per shot per AI.  The likelihood is extremely sharply
peaked in $\eta_{iz}$, so profiling and marginalising are numerically
equivalent.  We profile: for each shot $i$ and AI $z$, find
\begin{equation}
  \hat\eta_{iz}
  = \arg\max_{\eta}\;
    \Bigl[
      \operatorname{logsumexp}_{k}\!\log\Lcal_i^{(z)}(\eta,\phi_k)
      + \log p(\eta)
    \Bigr],
  \label{eq:eta_profile}
\end{equation}
where $\log\Lcal_i^{(z)}$ is the single-AI, single-shot conditional logL
using \emph{both} ports jointly:
\begin{equation}
  \log\Lcal_i^{(z)}(\eta,\phi_k)
  = \sum_{s \in \{g,e\}} \sum_b
    \Bigl[
      n_{izsb} \log \lambda_{izsb}(\phi_k;\eta)
      - \lambda_{izsb}(\phi_k;\eta)
    \Bigr],
\end{equation}
and $\log p(\eta)$ is a Gaussian prior (Section~\ref{sec:prior}).

Crucially, this profiling objective already marginalises over $\phi_i$ via
logsumexp, using the same phase grid as the signal fit.  The resulting
$\hat\eta_{iz}$ is then fixed for all subsequent evaluations of
$\log\Lcal(\beta)$.

\paragraph{Why $\hat\eta$ is approximately independent of $\beta$.}
The cloud parameters $\eta_{iz}$ control the spatial distribution of atoms
across pixels; $\beta$ controls the port fractions via the fringe.  These two
effects are largely orthogonal, so $\hat\eta_{iz}(\beta)$ changes negligibly
as $\beta$ varies.  This justifies computing $\hat\eta_{iz}$ once (at
arbitrary $\beta$, e.g.\ $\beta = 0$) and treating it as fixed during the
outer $\beta$ optimisation.

\subsection{Prior on Cloud Parameters}
\label{sec:prior}

The Gaussian log-prior on $\eta_{iz}$ is
\begin{align}
  \log p(\eta)
  &= -\frac{\mu_{x_0}^2}{2\sigma_{\mu,\mathrm{pos}}^2}
     -\frac{\mu_{y_0}^2}{2\sigma_{\mu,\mathrm{pos}}^2}
     -\frac{\mu_{v_x}^2}{2\sigma_{\mu,\mathrm{vel}}^2}
     -\frac{\mu_{v_y}^2}{2\sigma_{\mu,\mathrm{vel}}^2} \notag\\
  &\quad
     -\frac{(\sigma_x - \bar\sigma_\mathrm{pos})^2}{2\sigma_{\sigma,\mathrm{pos}}^2}
     -\frac{(\sigma_y - \bar\sigma_\mathrm{pos})^2}{2\sigma_{\sigma,\mathrm{pos}}^2}
     -\frac{(\sigma_{v_x} - \bar\sigma_\mathrm{vel})^2}{2\sigma_{\sigma,\mathrm{vel}}^2}
     -\frac{(\sigma_{v_y} - \bar\sigma_\mathrm{vel})^2}{2\sigma_{\sigma,\mathrm{vel}}^2},
  \label{eq:prior}
\end{align}
where the hyperparameters $(\sigma_{\mu,\mathrm{pos}},\sigma_{\mu,\mathrm{vel}},
\bar\sigma_\mathrm{pos},\bar\sigma_\mathrm{vel},\sigma_{\sigma,\mathrm{pos}},
\sigma_{\sigma,\mathrm{vel}})$ are parsed from the dataset name and can be
overridden by command-line arguments.

The optimisation uses a log-sigma parameterisation of the cloud widths:
$\tilde\sigma_x = \log\sigma_x$ etc., so that positivity of widths is enforced
by the domain of $\R$ without explicit bounds.  Bounded L-BFGS-B (rather
than Nelder--Mead) is used, with bounds on the log-widths derived from the
physical constraints of the image field of view.

\paragraph{Initialisation.}
Two strategies are supported:
\begin{itemize}
  \item \emph{Oracle init}: $\hat\eta_{iz}^{(0)} =$ true cloud parameters
        read from the H5 file.  Used for validation; referred to as
        \texttt{--init-from-true}.
  \item \emph{Prior-mean init}: $\hat\eta_{iz}^{(0)} = (0,0,0,0,\bar\sigma_\mathrm{pos},
        \bar\sigma_\mathrm{pos},\bar\sigma_\mathrm{vel},\bar\sigma_\mathrm{vel})$.
\end{itemize}

%%============================================================
\section{Alternative: Moment-Based Cloud Marginalisation}
\label{sec:moments}
%%============================================================

\subsection{Motivation}

The profiling strategy of Section~\ref{sec:nuisances} finds a single
$\hat\eta_{iz}$ per shot and fixes it.  This is valid when the likelihood is
sharply peaked in $\eta$, but empirical multi-start optimisation
(Section~\ref{sec:impl}) reveals that the likelihood surface has extended
\emph{ridges} in $\eta$-space, arising from ballistic degeneracies
(Section~\ref{sec:ridges}).  Fixing $\hat\eta$ at one point on a ridge
introduces a systematic error in the marginalised signal posterior.

This section develops an alternative that replaces the per-pixel likelihood
over all $N_b$ bins with a likelihood over a small set of image
\emph{moments}, enabling an analytic posterior $p(\eta|m)$ that concentrates
mass on the ridge and can be used as a proposal or prior for a proper
marginalisation over $\eta$.

\subsection{Ballistic Degeneracies and Ridge Structure}
\label{sec:ridges}

Under free-flight propagation~\eqref{eq:ballistic} the final-plane Gaussian
parameters are related to the initial ones by
\begin{align}
  \mu_{xf}   &= \mu_{x_0} + T_\mathrm{det}\,\mu_{v_x},
  \label{eq:mu_ballistic}\\
  \sigma^2_{xf} &= \sigma^2_{x_0} + T_\mathrm{det}^2\,\sigma^2_{v_x},
  \label{eq:sig_ballistic}
\end{align}
and identically for $y$.  The cloud image is predominantly determined by the
four \emph{final} moments $(\mu_{xf}, \mu_{yf}, \sigma_{xf}, \sigma_{yf})$.
Each of~\eqref{eq:mu_ballistic}--\eqref{eq:sig_ballistic} is a constraint
relating two initial parameters to one observable, leaving a one-dimensional
family of initial conditions consistent with any given image.  In
$(\mu_{x_0}, \mu_{v_x})$ space this is the line
$\mu_{x_0} + T_\mathrm{det}\,\mu_{v_x} = \mu_{xf}$; in
$(\sigma_{x_0},\sigma_{v_x})$ space it is the ellipse
$\sigma^2_{x_0} + T_\mathrm{det}^2\sigma^2_{v_x} = \sigma^2_{xf}$.  These
are precisely the ridges observed empirically.

\subsection{Summary Statistics and the Key Decomposition}

\paragraph{Total-count image.}
Define the per-pixel sum over output states
\begin{equation}
  m_b^{(iz)} = n_{izgb} + n_{izeb},
  \label{eq:mtot}
\end{equation}
and define the scalar image moments
\begin{align}
  N_z     &= \textstyle\sum_b m_b^{(iz)},
  \label{eq:Nz}\\
  \hat\mu_{xf} &= N_z^{-1}\textstyle\sum_b x_b\,m_b^{(iz)},
  \label{eq:muxf}\\
  \hat\sigma^2_{xf} &= N_z^{-1}\textstyle\sum_b (x_b - \hat\mu_{xf})^2\,m_b^{(iz)},
  \label{eq:sigxf}
\end{align}
and analogously $\hat\mu_{yf}$, $\hat\sigma^2_{yf}$.
We collect these into $m = (N_z, \hat\mu_{xf}, \hat\sigma^2_{xf}, \hat\mu_{yf},
\hat\sigma^2_{yf})$.

\paragraph{Science parameters do not enter $m$.}
The expected total-count rate at pixel $b$ is
\begin{equation}
  \E[m_b^{(iz)}]
  = L_{iz}\bigl[(A_{zgb}+A_{zeb})
    + (C^c_{zgb}+C^c_{zeb})\cos\Phi
    + (C^s_{zgb}+C^s_{zeb})\sin\Phi\bigr].
  \label{eq:Emtot}
\end{equation}
In a balanced two-path interferometer the output ports are complementary:
atoms transmitted into the ground port are not transmitted into the excited port,
so $A_{zgb}+A_{zeb} \approx \text{const}$ and
\begin{equation}
  C^c_{zgb}+C^c_{zeb} \approx 0,\qquad C^s_{zgb}+C^s_{zeb} \approx 0.
  \label{eq:complementarity}
\end{equation}
Therefore $\E[m_b^{(iz)}] \approx L_{iz}(A_{zgb}+A_{zeb})$, independent of
the total AI phase $\Phi$ and hence independent of both the laser phase
$\phi_i$ and the science signal $\beta = (A_s, A_c)$.  Consequently,
\begin{equation}
  \boxed{p(m \mid \eta,\,\beta,\,\phi_i) = p(m \mid \eta)}
  \label{eq:marg_key}
\end{equation}
exactly to the extent that~\eqref{eq:complementarity} holds.  The moments
are a sufficient statistic for $\eta$ that carries \emph{zero} information
about $\beta$.  This decomposition is not an approximation: the science signal
$\delta\phi_i(\beta)$ and the laser phase $\phi_i$ both enter the image
exclusively through the interference terms $C^c\cos\Phi + C^s\sin\Phi$,
which cancel in the port sum.

\subsection{Moment Likelihood $p(m \mid \eta)$}

Under~\eqref{eq:complementarity} the total-count image is a 2D Poisson
process with pixel rates $\lambda_b = L_{iz}\,A^{tot}_{zb}(\eta)$, where
$A^{tot}_{zb} = A_{zgb}+A_{zeb}$.  For the cloud model of
Section~\ref{sec:nuisances}, $A^{tot}_{zb}(\eta)$ is proportional to the
probability that a Gaussian-distributed atom lands in pixel $b$ at
detection time.

By the central limit theorem for large $N_z$, the sample moments converge to
\begin{align}
  \hat\mu_{xf} \mid \eta
  &\;\sim\; \mathcal{N}\!\Bigl(\mu_{xf}(\eta),\;\frac{\sigma^2_{xf}(\eta)}{N_z}\Bigr),
  \label{eq:mu_lik}\\
  \hat\sigma^2_{xf} \mid \eta
  &\;\sim\; \mathcal{N}\!\Bigl(\sigma^2_{xf}(\eta),\;\frac{2\,\sigma^4_{xf}(\eta)}{N_z}\Bigr),
  \label{eq:sig_lik}
\end{align}
and identically for $y$.  All four are independent.  The forward model
is~\eqref{eq:mu_ballistic}--\eqref{eq:sig_ballistic} with $N_z$ known from
the observed total count and $\sigma^2_{xf}$ estimated from the data
(replacing the true value in the noise variance by its estimate introduces
a negligible bias at large $N_z$).

\paragraph{Assumption: PSMAP undetectable in image moments.}
Equation~\eqref{eq:complementarity} is exact when port probabilities sum
to one.  In practice, paths are pruned (Section~2.2), so the fringe terms
do not cancel perfectly.  The residual fringe contrast is
$\epsilon_\mathrm{fringe} = |C^c_{zb}+C^c_{z(e)b}|/(A^{tot}_{zb}) \ll 1$.
Furthermore, the spatially varying phase $\Delta\phi_p(\q_0)$ (the PSMAP
contribution) averages over the cloud and partially cancels in the moment
integrals.  We assume this residual is negligible compared with the Poisson
uncertainty $1/\sqrt{N_z}$ on the moments, i.e.\ the PSMAP fringe structure
is undetectable in the image given the available atom number.
Under this assumption~\eqref{eq:marg_key} holds exactly and the
moment-based approach is not an approximation but an alternative sufficient
statistic.

\subsection{Analytic Posterior $p(\eta \mid m)$}

\subsubsection{Mean Parameters}

Consider the $x$-direction means $\theta_\mu = (\mu_{x_0}, \mu_{v_x})^\top$.
The observation model~\eqref{eq:mu_lik} combined with the
prior~\eqref{eq:prior} takes the form of a linear Gaussian model:
\begin{equation}
  \hat\mu_{xf} = \mathbf{h}^\top \theta_\mu + \epsilon,\qquad
  \mathbf{h} = \begin{pmatrix}1 \\ T_\mathrm{det}\end{pmatrix},\qquad
  \epsilon \sim \mathcal{N}(0, R_x),\quad
  R_x = \frac{\hat\sigma^2_{xf}}{N_z},
  \label{eq:linmod}
\end{equation}
with Gaussian prior $\theta_\mu \sim \mathcal{N}(\mathbf{0}, \mathbf{P}_0)$,
$\mathbf{P}_0 = \mathrm{diag}(\sigma^2_{\mu,\mathrm{pos}},\sigma^2_{\mu,\mathrm{vel}})$.
The posterior is analytic (standard Bayesian linear regression / Kalman update):
\begin{align}
  \mathbf{K}    &= \mathbf{P}_0\,\mathbf{h}
                   \bigl(\mathbf{h}^\top\mathbf{P}_0\,\mathbf{h} + R_x\bigr)^{-1},
  \label{eq:Kgain}\\
  \boldsymbol{\mu}_\mathrm{post} &= \mathbf{K}\,\hat\mu_{xf},
  \label{eq:mu_post}\\
  \mathbf{P}_\mathrm{post} &= (\mathbf{I} - \mathbf{K}\mathbf{h}^\top)\,\mathbf{P}_0.
  \label{eq:P_post}
\end{align}
The posterior $\theta_\mu \mid \hat\mu_{xf} \sim \mathcal{N}(\boldsymbol{\mu}_\mathrm{post},
\mathbf{P}_\mathrm{post})$ is a 2D Gaussian with $\mathbf{P}_\mathrm{post}$ having
one very small eigenvalue (the direction pinned by the constraint
$\mu_{x_0}+T_\mathrm{det}\mu_{v_x}=\hat\mu_{xf}$) and one larger eigenvalue
(the unconstrained direction along the ridge).
An identical update applies to $(\mu_{y_0}, \mu_{v_y})$ using $\hat\mu_{yf}$.

\subsubsection{Width Parameters}

Define the squared widths $s_x = \sigma^2_{x_0}$, $s_{vx} = \sigma^2_{v_x}$.
The observation model~\eqref{eq:sig_lik} becomes
\begin{equation}
  \hat\sigma^2_{xf} = s_x + T_\mathrm{det}^2\,s_{vx} + \epsilon_\sigma,\qquad
  \epsilon_\sigma \sim \mathcal{N}(0, R_\sigma),\quad
  R_\sigma = \frac{2\,\hat\sigma^4_{xf}}{N_z}.
  \label{eq:linmod_sig}
\end{equation}
This is linear in $(s_x, s_{vx})$, so the Kalman update~\eqref{eq:Kgain}--\eqref{eq:P_post}
applies with $\mathbf{h} = (1,\, T_\mathrm{det}^2)^\top$ and prior variance
on $(s_x, s_{vx})$ derived from~\eqref{eq:prior} by the delta method.
The non-negativity constraints $s_x, s_{vx} \ge 0$ truncate the posterior;
the truncation is negligible when the prior is well away from zero.

\paragraph{Dominant-velocity-spread regime.}
In the MAGIS-100 setup ($T_\mathrm{det}=3.8$\,s, $\sigma_{v_x}\sim 100$--$309\,\mu$m/s),
$T_\mathrm{det}\,\sigma_{v_x} \sim 380$--$1170\,\mu$m $\gg \sigma_{x_0}\sim 100\,\mu$m.
Expanding~\eqref{eq:sig_ballistic} to first order in $\sigma_{x_0}/\sigma_{xf}$:
\begin{equation}
  \sigma_{xf} \approx T_\mathrm{det}\,\sigma_{v_x}
  + \frac{\sigma^2_{x_0}}{2\,T_\mathrm{det}\,\sigma_{v_x}}.
  \label{eq:sig_expand}
\end{equation}
To leading order $\sigma_{v_x} \approx \hat\sigma_{xf}/T_\mathrm{det}$ is
directly observed, while $\sigma_{x_0}$ enters only as a small correction and
is poorly constrained by the image alone.  The posterior on $\sigma_{v_x}$
is therefore narrow ($\propto$ width of the Kalman update~\eqref{eq:P_post}),
while the posterior on $\sigma_{x_0}$ remains broad, reflecting its
unidentifiability from a single-epoch image.

\subsection{Marginalization via the Moment Posterior}

The moment posterior factorises as
\begin{equation}
  p(\eta \mid m)
  = p(\mu_{x_0},\mu_{v_x} \mid \hat\mu_{xf})
  \cdot p(\mu_{y_0},\mu_{v_y} \mid \hat\mu_{yf})
  \cdot p(\sigma^2_{x_0},\sigma^2_{v_x} \mid \hat\sigma^2_{xf})
  \cdot p(\sigma^2_{y_0},\sigma^2_{v_y} \mid \hat\sigma^2_{yf}),
  \label{eq:factorize}
\end{equation}
each factor being an analytic Gaussian (or truncated Gaussian) derived above.

The signal posterior is then
\begin{equation}
  p(\beta \mid \text{data})
  \;\propto\;
  \int d\eta\;
  p(\text{data} \mid \beta, \eta)\;
  p(\eta \mid m).
  \label{eq:marg_moment}
\end{equation}
This integral has two key advantages over the full marginalisation with
prior $p(\eta)$:

\begin{enumerate}
  \item \textbf{Constrained support.}  $p(\eta \mid m)$ concentrates mass on a
        4-dimensional manifold (the four ballistic ridges) rather than
        spanning the full 8D prior volume.  The effective sample size required
        for accurate Monte Carlo integration is thus dramatically reduced.

  \item \textbf{Analytic sampler.}  Samples $\eta_i \sim p(\eta \mid m)$ are
        free: draw from four independent 2D Gaussians.  No PSMAP evaluation
        is needed to generate starting points.
\end{enumerate}

\paragraph{Double-counting note.}
The moments $m$ are computed from the same data used in $p(\text{data}|\beta,\eta)$.
Double-counting is avoided because~\eqref{eq:marg_key} guarantees that $m$
carries no information about $\beta$: the moment channel is orthogonal to
the signal channel.  The only overlap is between $m$ and the cloud channel
of $p(\text{data}|\beta,\eta)$ — but this overlap vanishes by the
assumption that the PSMAP fringe is undetectable in the moments.  When this
assumption is not exact, one should use importance sampling to correct for
it (see Section~\ref{sec:is}).

\subsection{Analytical R\,=\,0 Limit: Dimension Reduction to a Point Estimate}
\label{sec:r0}

A practically important special case arises when we treat the sample moments
as \emph{exact} observations — i.e.\ set the measurement noise to zero
($R_x = R_\sigma = 0$ in~\eqref{eq:linmod}--\eqref{eq:linmod_sig}).
This is justified when $N_z \gg 1$ so that the CLT variance of the sample
moments is negligible relative to the prior spread.

\paragraph{R=0 Kalman posterior.}
Setting $R_x \to 0$, the Kalman gain in~\eqref{eq:Kgain} satisfies
\[
  \mathbf{K} \;\to\; \frac{\mathbf{P}_0\,\mathbf{h}}{\mathbf{h}^\top\mathbf{P}_0\,\mathbf{h}},
\]
and the posterior covariance $\mathbf{P}_\mathrm{post} \to \mathbf{P}_0 -
\mathbf{P}_0\,\mathbf{h}\,\mathbf{h}^\top\mathbf{P}_0 / (\mathbf{h}^\top\mathbf{P}_0\,\mathbf{h})$
is \emph{rank-deficient}: it has a zero eigenvalue in the $\mathbf{h}$ direction.
The posterior is a degenerate Gaussian — a delta function in the constraint direction
$\mathbf{h}^\top\theta = z$ (the ballistic manifold), and a 1D Gaussian in the
orthogonal (ridge) direction.

\paragraph{Dimension reduction.}
Each of the four 2D blocks
$(\mu_{x_0}, \mu_{v_x})$, $(\mu_{y_0}, \mu_{v_y})$,
$(\sigma^2_{x_0}, \sigma^2_{v_x})$, $(\sigma^2_{y_0}, \sigma^2_{v_y})$
has one degree of freedom eliminated by its R=0 constraint.  The full
8-dimensional nuisance space $\eta \in \mathbb{R}^8$ collapses to a
\textbf{4-dimensional} product of ridge manifolds:
\begin{equation}
  \{\eta : \mu_{x_0}+T\mu_{v_x} = \hat\mu_{xf}\}
  \;\times\;
  \{\eta : \mu_{y_0}+T\mu_{v_y} = \hat\mu_{yf}\}
  \;\times\;
  \{\eta : \sigma^2_{x_0}+T^2\sigma^2_{v_x} = \hat\sigma^2_{xf}\}
  \;\times\;
  \{\eta : \sigma^2_{y_0}+T^2\sigma^2_{v_y} = \hat\sigma^2_{yf}\}.
  \label{eq:ridge4d}
\end{equation}
The marginal~\eqref{eq:marg_moment} thus reduces to a \emph{4D integral}
over these four ridges.

\paragraph{Point-estimate approximation (zero integration).}
If we further approximate the marginal by evaluating the likelihood at a
\emph{single} representative point per ridge — the posterior mean,
which is the minimum-prior-cost point on each constraint manifold — the
integral reduces to a scalar evaluation.  For the mean block with diagonal
prior $P_0 = \mathrm{diag}(\tau_\mathrm{pos}^2, \tau_\mathrm{vel}^2)$
and measurement $\hat\mu_{xf}$ (exact), this gives:
\begin{equation}
  \hat\mu_{x_0} =
    \frac{\tau_\mathrm{pos}^2}{\tau_\mathrm{pos}^2 + T^2\,\tau_\mathrm{vel}^2}
    \,\hat\mu_{xf},
  \qquad
  \hat\mu_{v_x} =
    \frac{T\,\tau_\mathrm{vel}^2}{\tau_\mathrm{pos}^2 + T^2\,\tau_\mathrm{vel}^2}
    \,\hat\mu_{xf},
  \label{eq:r0mean}
\end{equation}
and analogously for $y$ and the variance blocks.  No integration or Monte Carlo
is required.

\paragraph{Bias and validity of the point estimate.}
The true $\eta$ lies on the constraint manifold~\eqref{eq:ridge4d} but in
general not at the minimum-norm point~\eqref{eq:r0mean}.  Using the point
estimate therefore introduces a bias in the ACS evaluation unless the fringe
coefficients are constant along the ridge.  Section~\ref{sec:gh4d} describes
the unbiased 4D integral implemented in \texttt{python-scripts/moment\_inference.py}.

\subsection{Importance Sampling Correction and Information Diagnostic}
\label{sec:is}

If the PSMAP fringe is not fully cancelled in the moment integrals, using
$p(\eta|m)$ as a prior introduces a small bias.  This is corrected exactly
by importance sampling:

\begin{enumerate}
  \item Draw $M$ samples $\eta_j \sim q(\eta) = p(\eta \mid m_\alpha)$, where
        $m_\alpha$ is the moment posterior with covariances inflated by a
        factor $\alpha > 1$ (e.g.\ $\alpha = 4$) to ensure the proposal has
        heavier tails than the true posterior.

  \item Evaluate the unnormalised importance weights
        \begin{equation}
          w_j = \frac{p(\eta_j)}{q(\eta_j)}
              \;\propto\;
              \frac{p(\eta_j)}{p(\eta_j \mid m_\alpha)},
          \label{eq:isweights}
        \end{equation}
        which correct for using $q$ in place of the flat prior $p(\eta)$.

  \item Approximate the integral~\eqref{eq:marg_moment} by
        \begin{equation}
          p(\beta \mid \text{data})
          \;\propto\;
          \sum_{j=1}^M w_j\;p(\text{data} \mid \beta, \eta_j).
          \label{eq:isapprox}
        \end{equation}
\end{enumerate}

\paragraph{Information diagnostic.}
The effective sample size
\begin{equation}
  \mathrm{ESS} = \frac{(\sum_j w_j)^2}{\sum_j w_j^2}
  \label{eq:ess}
\end{equation}
measures how much the full pixel likelihood reshapes $p(\eta)$ relative to
$p(\eta|m)$.  If $\mathrm{ESS}/M \approx 1$, the moment posterior
$p(\eta|m)$ is essentially the correct posterior — the PSMAP fringe structure
adds no detectable information about $\eta$ beyond the cloud envelope moments,
confirming the assumption of Section~\ref{sec:moments}.  If
$\mathrm{ESS}/M \ll 1$, the PSMAP contributes additional cloud information
and full joint inference is required.

%%============================================================
\section{4D Gauss-Hermite Moment Inference}
\label{sec:gh4d}
%%============================================================

This section describes the complete, unbiased implementation of the R=0 ridge
integral introduced in Section~\ref{sec:r0}.  Rather than fixing $\eta$ at the
single minimum-norm ridge point, we integrate over all initial conditions
consistent with the observed final moments using Gauss-Hermite quadrature with
an analytic weight function derived from the R=0 posterior.  No per-pixel
spatial information is used; the sufficient statistic is reduced to five scalars
per shot per AI.

\subsection{Spatially-Integrated ACS}
\label{sec:acs_integrated}

Instead of the per-pixel ACS tensors $\{A_{zsb},\,C^c_{zsb},\,C^s_{zsb}\}$ of
Section~\ref{sec:qmc}, we work with their spatial sums,
\begin{equation}
  A_{zs}(\eta) = \sum_b A_{zsb}(\eta),\qquad
  C^c_{zs}(\eta) = \sum_b C^c_{zsb}(\eta),\qquad
  C^s_{zs}(\eta) = \sum_b C^s_{zsb}(\eta),
  \label{eq:acs_int}
\end{equation}
reducing the cloud-marginalized pixel ACS to three scalars per state per AI.
These are the expectations of the PSMAP amplitudes over the full cloud
Gaussian, without conditioning on final position:
\begin{align}
  A_{zs}(\eta) &= \E_{\q_0 \sim \mathcal{N}(\eta)}
    \!\Bigl[\textstyle\sum_{p:\,s_p=s}
      \bigl(a_{0,p}^2(\q_0)+a_{1,p}^2(\q_0)\bigr)\Bigr],
  \label{eq:Aint}\\[3pt]
  C^c_{zs}(\eta) &= \E_{\q_0 \sim \mathcal{N}(\eta)}
    \!\Bigl[\textstyle\sum_{p:\,s_p=s}
      \iota_p\,2\,a_{0,p}(\q_0)\,a_{1,p}(\q_0)\cos\Delta\phi_p(\q_0)\Bigr],\\[3pt]
  C^s_{zs}(\eta) &= \E_{\q_0 \sim \mathcal{N}(\eta)}
    \!\Bigl[-\textstyle\sum_{p:\,s_p=s}
      \iota_p\,2\,a_{0,p}(\q_0)\,a_{1,p}(\q_0)\sin\Delta\phi_p(\q_0)\Bigr].
\end{align}
The expected detected count for state $s$ at total phase $\Phi$ is then
\begin{equation}
  \lambda_{zs}(\Phi;\,\eta)
  = \Lambda_{iz}\,\bigl[A_{zs}(\eta)+C^c_{zs}(\eta)\cos\Phi
    +C^s_{zs}(\eta)\sin\Phi\bigr],
  \label{eq:lambda_int}
\end{equation}
where $\Lambda_{iz}$ is an atom-count normalisation factor defined below.

\paragraph{Monte Carlo evaluation.}
Let $n_\mathrm{qmc}$ standard-normal deviates $\mathbf{z}_j\!\in\!\mathbb{R}^4$
be drawn once.  Given $\eta$, the cloud samples are
$\q_{0,j}=\mu(\eta)+\mathrm{diag}(\sigma(\eta))\,\mathbf{z}_j$.  Evaluate the
PSMAP at all $n_\mathrm{qmc}$ points in a single batched GPU call, then
average.  Because the samples are not conditioned on final position, the
estimator variance is $\mathcal{O}(1/n_\mathrm{qmc})$ without the
$1/n_\mathrm{occupied}$ binning penalty of the per-pixel estimator.  The default
$n_\mathrm{qmc}=2048$ is sufficient because the required ACS accuracy is set by
per-shot signal-to-noise, not per-pixel count statistics.

\subsection{Port-Fraction Likelihood}
\label{sec:port_frac}

The observations per shot per AI are reduced to two scalars: total
ground-state count $n_{izg}=\sum_b n_{izgb}$ and total excited-state count
$n_{ize}=\sum_b n_{izeb}$, with $N_{iz}=n_{izg}+n_{ize}$.  Define
\begin{equation}
  A^{\rm tot}_z=A_{zg}+A_{ze},\quad
  C^{c,\rm tot}_z=C^c_{zg}+C^c_{ze},\quad
  C^{s,\rm tot}_z=C^s_{zg}+C^s_{ze}.
\end{equation}
Because path-pruning violates complementarity ($C^{c,\rm tot}_z\neq0$ in
general), the total expected count at phase $\Phi$ is $\Phi$-dependent:
\begin{equation}
  \lambda^{\rm tot}_{iz}(\Phi;\,\eta)
  = \Lambda_{iz}\bigl[A^{\rm tot}_z+C^{c,\rm tot}_z\cos\Phi
    +C^{s,\rm tot}_z\sin\Phi\bigr].
  \label{eq:lambda_tot}
\end{equation}
We normalise by choosing $\Lambda_{iz}(\Phi;\eta)$ so that the predicted total
matches the observed total at every phase:
\begin{equation}
  \Lambda_{iz}(\Phi;\,\eta)
  = \frac{N_{iz}}{A^{\rm tot}_z(\eta)+C^{c,\rm tot}_z(\eta)\cos\Phi
    +C^{s,\rm tot}_z(\eta)\sin\Phi}.
  \label{eq:Lphi}
\end{equation}
This $\Phi$-dependent normalisation eliminates the total-count channel,
leaving the log-likelihood as a function of the port fraction only:
\begin{align}
  \hat\lambda_{izg}(\Phi;\,\eta)
  &= N_{iz}\,\frac{A_{zg}+C^c_{zg}\cos\Phi+C^s_{zg}\sin\Phi}
    {A^{\rm tot}_z+C^{c,\rm tot}_z\cos\Phi+C^{s,\rm tot}_z\sin\Phi},
  \label{eq:lam_g_normed}\\
  \hat\lambda_{ize}(\Phi;\,\eta)
  &= N_{iz}-\hat\lambda_{izg}(\Phi;\,\eta).
\end{align}
The Poisson log-likelihood for AI $z$ at phase $\Phi$ is
\begin{equation}
  \ell_{iz}(\Phi;\,\eta)
  = n_{izg}\log\hat\lambda_{izg}
    -\hat\lambda_{izg}
    +n_{ize}\log\hat\lambda_{ize}
    -\hat\lambda_{ize},
  \label{eq:ll_ai}
\end{equation}
which depends only on the port fraction $p_g(\Phi;\eta)=\hat\lambda_{izg}/N_{iz}$.
The combined Z0$+$Z100 conditional log-likelihood for shot $i$ at laser phase
$\phi$ is
\begin{equation}
  \ell_i(\phi;\,\eta_{iz_0},\eta_{iz_{100}},\beta)
  = \ell_{iz_0}(\phi;\,\eta_{iz_0})
  +\ell_{iz_{100}}(\phi+\delta\phi_i(\beta);\,\eta_{iz_{100}}),
  \label{eq:ll_combined}
\end{equation}
where $\delta\phi_i(\beta)=A_s\sin(2\pi f\,i)+A_c\cos(2\pi f\,i)$ is the
differential phase signal injected into Z100.

\subsection{1D Ridge Prior from the R=0 Kalman Update}
\label{sec:ridge_prior_gh}

Under the R=0 Kalman update (Section~\ref{sec:r0}), the posterior for the
mean block $(\mu_{x_0},\mu_{v_x})$ conditional on $\hat\mu_{xf}$ is a
degenerate Gaussian with a non-degenerate 1D marginal in the direction
parametrised by $\mu_{x_0}$.  Substituting the R=0 Kalman gain gives
\begin{equation}
  \mu_{x_0}\mid\hat\mu_{xf}
  \;\sim\;
  \mathcal{N}\!\bigl(\bar\mu_{x_0},\;\bar\sigma^2_{\mu_{x_0}}\bigr),
  \label{eq:ridge_prior_mu}
\end{equation}
\begin{equation}
  \bar\mu_{x_0}
  = \frac{\tau_\mathrm{pos}^2}{\tau_\mathrm{pos}^2+T^2\tau_\mathrm{vel}^2}
    \hat\mu_{xf},
  \qquad
  \bar\sigma^2_{\mu_{x_0}}
  = \!\left(\frac{1}{\tau_\mathrm{pos}^2}
    +\frac{1}{T^2\tau_\mathrm{vel}^2}\right)^{\!-1}.
  \label{eq:ridge_prior_mu_params}
\end{equation}
Once $\mu_{x_0}$ is fixed, the ridge constraint gives
$\mu_{v_x}=(\hat\mu_{xf}-\mu_{x_0})/T$.
An identical pair of equations applies to $(\mu_{y_0},\mu_{v_y})$.

For the variance block, a linearised prior on $s_{x_0}=\sigma^2_{x_0}$,
$s_{v_x}=\sigma^2_{v_x}$ and the R=0 Kalman update on the constraint
$s_{x_0}+T^2 s_{v_x}=\hat\sigma^2_{xf}$ give
\begin{equation}
  s_{x_0}\mid\hat\sigma^2_{xf}
  \;\sim\;
  \mathcal{N}\!\bigl(\bar s^{\,*}_{x_0},\;\bar\tau^2_{s_{x_0}}\bigr),
  \label{eq:ridge_prior_var}
\end{equation}
\begin{equation}
  \bar s^{\,*}_{x_0}
  = \bar\sigma_\mathrm{pos}^2
  + \bar\tau^2_{s_{x_0}}\,
    \frac{\hat\sigma^2_{xf}-\bar\sigma_\mathrm{pos}^2-T^2\bar\sigma_\mathrm{vel}^2}
         {\tau^2_{s_{x_0}}},
  \qquad
  \bar\tau^2_{s_{x_0}}
  = \!\left(\frac{1}{\tau^2_{s_{x_0}}}
    +\frac{1}{T^4\tau^2_{s_{v_x}}}\right)^{-1},
  \label{eq:ridge_prior_var_params}
\end{equation}
where $\tau_{s_{x_0}}=2\bar\sigma_\mathrm{pos}\sigma_{\sigma,\mathrm{pos}}$ is
the prior standard deviation on $s_{x_0}$ obtained by the delta method, and
$\bar\sigma_\mathrm{pos}$, $\sigma_{\sigma,\mathrm{pos}}$ are the prior mean
and standard deviation of $\sigma_{x_0}$.  Once $s_{x_0}$ is fixed,
$s_{v_x}=\max(\hat\sigma^2_{xf}-s_{x_0},\,0)/T^2$.  These four 1D priors are
the analytic Bayesian combination of the original 8D prior $p(\eta)$ with the
observed final moments, not a heuristic.

\subsection{4D Gauss-Hermite Quadrature Grid}
\label{sec:gh_grid}

The 4D nuisance integral after the R=0 reduction is
\begin{equation}
  \int_{\mathbb{R}^4}
    \ell_i\!\bigl(\phi;\,\tilde\eta(\xi),\beta\bigr)\;
    p_{\rm ridge}(\xi\mid m_i)\;d^4\xi,
  \label{eq:ridge_int}
\end{equation}
where $\xi=(\mu_{x_0},\mu_{y_0},s_{x_0},s_{y_0})$ collects the four free
parameters, $\tilde\eta(\xi)$ reconstructs the full 8-vector via the ridge
constraints, and $p_{\rm ridge}=\prod_{d=1}^4\mathcal{N}(\bar\xi_d,\bar\tau_d^2)$
is the product of four independent 1D Gaussians from
Section~\ref{sec:ridge_prior_gh}.

Since $p_{\rm ridge}$ is a product of Gaussians, the integral is approximated
by a 4D product Gauss-Hermite (GH) rule.  With $n_{\rm gh}$ nodes per
dimension, the rule uses $n_{\rm gh}^4$ points and is exact for polynomials
of degree $\le2n_{\rm gh}-1$ in each dimension.

\paragraph{Construction.}
Let $\{t_k,w_k\}_{k=1}^{n_{\rm gh}}$ be the standard Hermite nodes and
weights ($\sum_k w_k=\sqrt{\pi}$).  For dimension $d$ with Gaussian
$\mathcal{N}(\bar\xi_d,\bar\tau_d^2)$, the physical nodes and normalised
log-weights are
\begin{equation}
  \xi_{d,k}=\bar\xi_d+\sqrt{2}\,\bar\tau_d\,t_k,
  \qquad
  \log\tilde w_{d,k}=\log w_k-\tfrac{1}{2}\log\pi,
  \label{eq:gh_nodes}
\end{equation}
with $\sum_k\tilde w_{d,k}=1$.  The 4D product rule has node
$\xi^{(j)}=(\xi_{1,k_1},\xi_{2,k_2},\xi_{3,k_3},\xi_{4,k_4})$ and
\begin{equation}
  \log\tilde w^{(j)}
  =\sum_{d=1}^4\log\tilde w_{d,k_d}.
  \label{eq:gh_logw}
\end{equation}
Variance nodes $s_{x_0}$ are clipped to $[0,\hat\sigma^2_{xf}]$ before
reconstructing $\tilde\eta$:
\begin{equation}
  \mu_{v_x}=(\hat\mu_{xf}-\mu_{x_0})/T,\quad
  \sigma_{v_x}=\sqrt{\max(\hat\sigma^2_{xf}-s_{x_0},0)}/T,
  \label{eq:eta_reconstruct}
\end{equation}
and analogously for the $y$ block.

\subsection{Per-Shot Log-Marginal Likelihood}
\label{sec:shot_logmarg}

For each GH node $j=1,\ldots,n_{\rm gh}^4$, the integrated ACS
$\{A^{(j)}_{zs},C^{c(j)}_{zs},C^{s(j)}_{zs}\}$ is evaluated once via the
batched Monte Carlo estimator of Section~\ref{sec:acs_integrated} at
$\tilde\eta^{(j)}$, and cached.  All subsequent evaluations at different
$(\phi,\beta)$ reuse these precomputed values.

The GH-marginalised log-likelihood for AI $z$ at phase $\Phi$ is
\begin{equation}
  \log f_{iz}(\Phi)
  =\operatorname{logsumexp}_{j=1}^{n_{\rm gh}^4}
    \Bigl[\log\tilde w^{(j)}+\ell_{iz}\!\bigl(\Phi;\;\tilde\eta^{(j)}\bigr)\Bigr].
  \label{eq:logf_iz}
\end{equation}
The phase-marginalised log-likelihood for shot $i$ is obtained by combining
the two AIs and marginalising $\phi_i$ over a uniform grid
$\phi_k=2\pi k/K$, $k=0,\ldots,K-1$:
\begin{equation}
  \log\mathcal{L}_i(\beta)
  =\operatorname{logsumexp}_{k=0}^{K-1}
    \Bigl[\log f_{iz_0}(\phi_k)
    +\log f_{iz_{100}}(\phi_k+\delta\phi_i(\beta))\Bigr]-\log K.
  \label{eq:logLi}
\end{equation}
The outer logsumexp marginalises over the unknown laser phase $\phi_i$ with a
uniform discrete prior.  The full log-likelihood sums over independent shots:
\begin{equation}
  \log\mathcal{L}(\beta)=\sum_{i=1}^{N_{\rm shots}}\log\mathcal{L}_i(\beta).
  \label{eq:logLfull_gh}
\end{equation}

\subsection{Signal Optimisation}
\label{sec:signal_opt_gh}

The MLE $\hat\beta=\arg\max_\beta\log\mathcal{L}(\beta)$ is found by
multi-start L-BFGS-B.  One start is fixed at the null $(0,0)$; the remaining
$n_{\rm starts}-1$ are drawn uniformly from the search domain
$[-A_{\rm max},A_{\rm max}]^2$.  The best solution across all starts is
retained.  With $N_{\rm shots}\ge100$ and signal amplitude
$A\gtrsim0.05$\,rad, all starts converge to the same optimum, confirming
unimodality at this regime.

\subsection{Complete Algorithm}
\label{sec:gh_algorithm}

\begin{enumerate}

\item \textbf{Data reduction} (once per shot per AI).
  From the 2D image pair $(n_{izgb},n_{izeb})$, compute port sums
  $n_{izg}$, $n_{ize}$, $N_{iz}$, and spatial moments
  $(\hat\mu_{xf},\hat\mu_{yf},\hat\sigma^2_{xf},\hat\sigma^2_{yf})$
  from the combined image $m_b=n_{izgb}+n_{izeb}$.

\item \textbf{Ridge prior} (once per shot per AI).
  Compute the four 1D Gaussian parameters
  $\{(\bar\mu_{x_0},\bar\sigma_{\mu_{x_0}}),(\bar\mu_{y_0},\bar\sigma_{\mu_{y_0}}),
  (\bar s^{\,*}_{x_0},\bar\tau_{s_{x_0}}),(\bar s^{\,*}_{y_0},\bar\tau_{s_{y_0}})\}$
  via \eqref{eq:ridge_prior_mu_params} and \eqref{eq:ridge_prior_var_params}.

\item \textbf{GH grid} (once per shot per AI).
  Form the $n_{\rm gh}^4$ product-rule nodes $\{\xi^{(j)}\}$ and
  log-weights $\{\log\tilde w^{(j)}\}$ via \eqref{eq:gh_nodes}--\eqref{eq:gh_logw}.
  Clip variance nodes and reconstruct $\tilde\eta^{(j)}$ via \eqref{eq:eta_reconstruct}.

\item \textbf{Integrated ACS precomputation} (once per GH node per shot per AI).
  For each $\tilde\eta^{(j)}$, draw $n_{\rm qmc}$ Gaussian cloud samples and
  evaluate the batched PSMAP on GPU to obtain
  $(A^{(j)}_{zs},C^{c(j)}_{zs},C^{s(j)}_{zs})$ for $s\in\{g,e\}$.
  Store alongside $\log\tilde w^{(j)}$.

\item \textbf{Signal log-likelihood} (per $\beta$ evaluation).
  For each shot $i$ and trial $\beta$: compute $\delta\phi_i(\beta)$;
  evaluate $\log f_{iz}(\phi_k)$ for all $k$ via \eqref{eq:logf_iz};
  apply \eqref{eq:logLi}; sum shots via \eqref{eq:logLfull_gh}.

\item \textbf{Optimisation}.
  Maximise $\log\mathcal{L}(\beta)$ over $\beta\in[-A_{\rm max},A_{\rm max}]^2$
  with $n_{\rm starts}$ independent L-BFGS-B runs
  (Section~\ref{sec:signal_opt_gh}).

\end{enumerate}

\subsection{Default Hyperparameters and Computational Cost}
\label{sec:gh_cost}

\begin{table}[h]
\centering
\begin{tabular}{lll}
\toprule
Parameter & Symbol & Default \\
\midrule
Detection time                 & $T$ & $3.8$\,s \\
GH order per dimension         & $n_{\rm gh}$ & $4$ \\
Total GH nodes per shot per AI & $n_{\rm gh}^4$ & $256$ \\
QMC samples per GH node        & $n_{\rm qmc}$ & $2048$ \\
Phase grid points              & $K$ & $128$ \\
Optimiser starts               & $n_{\rm starts}$ & $8$ \\
Search domain half-width       & $A_{\rm max}$ & $0.2$\,rad \\
COM prior std (position)       & $\tau_{\rm pos}$ & $10\,\mu$m \\
COM prior std (velocity)       & $\tau_{\rm vel}$ & $10\,\mu$m/s \\
Width prior mean (position)    & $\bar\sigma_{\rm pos}$ & $100\,\mu$m \\
Width prior mean (velocity)    & $\bar\sigma_{\rm vel}$ & $100\,\mu$m/s \\
Width prior std (position)     & $\sigma_{\sigma,{\rm pos}}$ & $10\,\mu$m \\
Width prior std (velocity)     & $\sigma_{\sigma,{\rm vel}}$ & $10\,\mu$m/s \\
\bottomrule
\end{tabular}
\caption{Default hyperparameters for the 4D GH moment inference pipeline
(\texttt{python-scripts/moment\_inference.py}).}
\label{tab:gh_params}
\end{table}

\begin{table}[h]
\centering
\begin{tabular}{lll}
\toprule
Stage & Scope & Wall-time (A100) \\
\midrule
Load PSMAPs                       & Once               & ${\approx}25$\,s \\
Data reduction + ridge priors     & Per shot/AI        & ${<}1$\,ms \\
GH grid + ACS precompute          & Per shot/AI        & ${\approx}0.56$\,s \\
Total precompute (200 shots)      & Per run            & ${\approx}1.9$\,min \\
$\log\mathcal{L}(\beta)$ single eval & Per eval        & ${\approx}0.5$\,s \\
Optimisation ($n_{\rm starts}=8$) & Per run           & ${\approx}95$\,s \\
\textbf{Total (200 shots)}        & \textbf{Per run}  & $\mathbf{{\approx}10}$\,\textbf{min} \\
\bottomrule
\end{tabular}
\caption{Wall-times for $N_{\rm shots}=200$, $n_{\rm gh}=4$, $n_{\rm qmc}=2048$, $K=128$.
The precompute stage is dominated by $200\times2\times256=102\,400$ batched PSMAP calls.}
\label{tab:gh_cost}
\end{table}

\subsection{Validation}
\label{sec:gh_validation}

Running \texttt{moment\_inference.py} on the $N_{\rm shots}=200$, $N_{iz}=10^6$
dataset (true signal $A=0.10$\,rad, $f=0.30$\,cycles/shot, $\varphi=0.50$\,rad,
giving $(A_s,A_c)=(0.0878,0.0479)$\,rad) gives:
\begin{itemize}
  \item $\log\mathcal{L}(\hat\beta)-\log\mathcal{L}(0,0)=+1.79\times10^5$\,nats,
    strongly detecting the signal.
  \item $\hat A_s=0.091$\,rad, $\hat A_c=0.051$\,rad
    (recovered amplitude $\hat A=0.104$\,rad vs.\ true $0.100$\,rad).
  \item All 8 optimisation starts converge to the same value
    ($\log\mathcal{L}\approx4.098\times10^9$), confirming unimodality.
\end{itemize}

%%============================================================
\section{8D Prior Marginalisation (Moment-Free Baseline)}
\label{sec:prior8d}
%%============================================================

This section describes a simplified variant of the inference pipeline that
\emph{ignores} the observed final-position moments entirely.  Rather than
collapsing the nuisance integral to a 4D ridge via the R=0 Kalman update
(Section~\ref{sec:gh4d}), the cloud parameters $\eta$ are marginalised
directly against the 8D prior $p_0(\eta)$.  The result is a strictly
\emph{less informative} likelihood that serves as a baseline: comparing it
with the moment-conditioned result quantifies how much information the spatial
image data contribute beyond the total port counts.

\subsection{Reduction to Port Counts Only}

For each shot $i$ and each AI $z$, the only data used are the total port
counts $n_{izg} = \sum_b n_{izgb}$ and $n_{ize} = \sum_b n_{izeb}$.  No
spatial moments are extracted; the full image is reduced to two scalars.

The integrated ACS model~\eqref{eq:lam_g_normed}--\eqref{eq:ll_ai} and the
port-fraction log-likelihood~\eqref{eq:ll_combined} are used unchanged.

\subsection{Marginalisation over the 8D Prior}
\label{sec:prior8d_marg}

The log-marginal for shot $i$ at laser phase $\phi$ and signal $\beta$ is
\begin{equation}
  \log\mathcal{L}_i^{\rm prior}(\beta)
  = \operatorname{logsumexp}_{k=0}^{K-1}
    \Bigl[
      \log F_{iz_0}(\phi_k)
      + \log F_{iz_{100}}(\phi_k + \delta\phi_i(\beta))
    \Bigr] - \log K,
  \label{eq:logLi_prior}
\end{equation}
where
\begin{equation}
  \log F_{iz}(\Phi)
  = \operatorname{logsumexp}_{j=1}^{n_\eta}
    \Bigl[\log\tilde w_j + \ell_{iz}\!\bigl(\Phi;\;\eta^{(j)}\bigr)\Bigr],
  \qquad
  \log\tilde w_j = -\log n_\eta \;\;\forall j,
  \label{eq:logF_iz_prior}
\end{equation}
and $\{\eta^{(j)}\}_{j=1}^{n_\eta}$ are fixed samples from the prior
$p_0(\eta)$.  The uniform weights $\tilde w_j = 1/n_\eta$ reflect that
the samples are drawn directly from $p_0(\eta)$, so no importance
correction is needed.

\subsection{Sobol QMC Sampling of the Prior}
\label{sec:prior8d_qmc}

The prior factorises over the 8 cloud parameters
$\eta = (\mu_{x_0}, \mu_{y_0}, \mu_{v_x}, \mu_{v_y}, \sigma_{x_0},
\sigma_{y_0}, \sigma_{v_x}, \sigma_{v_y})$:
\begin{align}
  \mu_{x_0},\mu_{y_0} &\sim \mathcal{N}(0,\,\tau_\mathrm{pos}^2), &
  \mu_{v_x},\mu_{v_y} &\sim \mathcal{N}(0,\,\tau_\mathrm{vel}^2), \nonumber\\
  \sigma_{x_0},\sigma_{y_0} &\sim \mathcal{N}(\bar\sigma_\mathrm{pos},\,\sigma_{\sigma,\mathrm{pos}}^2), &
  \sigma_{v_x},\sigma_{v_y} &\sim \mathcal{N}(\bar\sigma_\mathrm{vel},\,\sigma_{\sigma,\mathrm{vel}}^2).
  \label{eq:prior8d}
\end{align}
Standard deviations are clipped to a small positive floor after sampling.

Samples are drawn using a scrambled Sobol sequence in $[0,1)^8$ transformed
through the normal quantile function.  Sobol sequences achieve a
discrepancy of $\mathcal{O}((\log n_\eta)^8 / n_\eta)$ in 8D, compared to
$\mathcal{O}(1/\sqrt{n_\eta})$ for plain Monte Carlo, yielding significantly
better coverage of the prior at the same sample budget.  The default
$n_\eta = 512$ gives a comparable total evaluation count to the $n_{\rm gh}^4
= 256$ nodes of the 4D GH rule.

\paragraph{Why not 8D Gauss-Hermite.}
A product GH rule with $n_{\rm gh}$ nodes per dimension requires
$n_{\rm gh}^8$ evaluations.  Even $n_{\rm gh}=2$ gives $2^8=256$ nodes
(only three polynomial degrees per axis — far too coarse for the nonlinear
ACS integrand).  $n_{\rm gh}=4$ would require $4^8=65\,536$ nodes, a 256-fold
increase over the 4D case, making it computationally infeasible.  QMC with
$n_\eta = 512$ achieves comparable or better accuracy at the same cost as the
4D moment-conditioned method.

\subsection{Computational Cost}

Because the prior $p_0(\eta)$ does not depend on any shot-specific data,
the integrated ACS need be evaluated at the $n_\eta$ prior samples
\emph{only once} per AI — not once per shot.  The precompute is therefore
$\mathcal{O}(n_\eta)$ instead of $\mathcal{O}(N_{\rm shots}\times n_\eta)$,
reducing the precompute from $\sim\!2$\,min to a few seconds.  Only the
per-shot port counts $(n_{izg}, n_{ize})$ are read in a fast loop.

\begin{table}[h]
\centering
\begin{tabular}{lll}
\toprule
Stage & Scope & Wall-time (A100) \\
\midrule
Load PSMAPs                                 & Once           & ${\approx}25$\,s \\
Sample $n_\eta$ prior points (Sobol)        & Once           & ${<}0.1$\,s \\
ACS precompute ($n_\eta$ pts, 2 AIs)        & Once           & ${\approx}5$\,s \\
Read port counts (200 shots)                & Once           & ${<}1$\,s \\
$\log\mathcal{L}^{\rm prior}(\beta)$ eval   & Per eval       & ${\approx}0.5$\,s \\
Optimisation ($n_{\rm starts}=8$)          & Per run        & ${\approx}95$\,s \\
\textbf{Total (200 shots)}                  & \textbf{Per run} & $\mathbf{{\approx}2}$\,\textbf{min} \\
\bottomrule
\end{tabular}
\caption{Wall-times for \texttt{python-scripts/prior\_inference.py},
$N_{\rm shots}=200$, $n_\eta=512$, $K=128$.  The precompute saving over the
moment-conditioned method (Section~\ref{sec:gh4d}) arises because the prior
ACS is shot-independent.}
\label{tab:prior_cost}
\end{table}

\subsection{Information Content of Final Moments}
\label{sec:moment_info}

Let $\hat\beta^{\rm mom}$ and $\hat\beta^{\rm prior}$ denote the MLEs from
the moment-conditioned and prior-only pipelines respectively.  Three
complementary diagnostics quantify the value of the spatial moment data:

\begin{enumerate}
  \item \textbf{Log-likelihood gain.}
    The quantity
    \begin{equation}
      \Delta\log\mathcal{L}
      = \log\mathcal{L}^{\rm mom}(\hat\beta^{\rm mom})
        - \log\mathcal{L}^{\rm prior}(\hat\beta^{\rm prior})
    \end{equation}
    is always non-negative (both methods use the same port-count data; the
    moment-conditioned method uses strictly more information).  A large
    $\Delta\log\mathcal{L}$ means the spatial moments substantially sharpen
    the nuisance posterior, aiding signal recovery.

  \item \textbf{Recovery accuracy.}
    Compare $|\hat\beta^{\rm mom} - \beta_{\rm true}|$ vs.\
    $|\hat\beta^{\rm prior} - \beta_{\rm true}|$.  Bias or increased variance
    in the prior-only estimate indicates that moment conditioning improves
    the identifiability of the signal.

  \item \textbf{Signal-to-noise ratio.}
    On repeated datasets (varying shot noise and cloud realisations), the
    ratio of signal-to-null log-likelihood ratios
    $\log\mathcal{L}(\hat\beta)/\log\mathcal{L}(0,0)$ from the two
    pipelines gives the effective SNR gain from moment conditioning.
\end{enumerate}

The prior-only method therefore serves as both a practical fallback (for
settings where images are not recorded or not trusted) and as a diagnostic
probe of how much spatial information is present in the cloud images.

%%============================================================
\section{Full Marginal Log-Likelihood}
%%============================================================

\subsection{Rotation Trick for AI $z=100$}

For AI $z=100$, the effective phase at grid point $k$ in shot $i$ is
$\phi_k + \delta\phi_i(\beta)$.  Rather than recomputing the ACS for each
$(\delta\phi_i, k)$, we apply a rotation to the precomputed ACS:
\begin{align}
  C^{c,\mathrm{eff}}_{zsb}(\beta)
  &= C^c_{zsb}\cos(\delta\phi_i) + C^s_{zsb}\sin(\delta\phi_i), \\
  C^{s,\mathrm{eff}}_{zsb}(\beta)
  &= -C^c_{zsb}\sin(\delta\phi_i) + C^s_{zsb}\cos(\delta\phi_i).
\end{align}
This follows from the angle-addition identity and means that the per-evaluation
cost for Z100 is a 2D rotation of the precomputed ACS tensors, not a new PSMAP
evaluation.

\subsection{Per-Shot Marginal Log-Likelihood}

Given precomputed ACS $\{A_{zsb}, C^c_{zsb}, C^s_{zsb}\}$ and scale factor
$L_{iz}$ at $\hat\eta_{iz}$, the conditional log-likelihood at phase grid
point $k$ is:
\begin{align}
  \log\Lcal_i(\beta,\phi_k)
  &= \sum_z \sum_{s \in \{g,e\}} \sum_b
     \Bigl[n_{izsb}\log\lambda_{izsb}(\phi_k)
           - \lambda_{izsb}(\phi_k)\Bigr],
  \label{eq:ll_phik}
\end{align}
where for Z0: $\lambda_{igsb}(\phi_k) = L_{i,0}\,[A_{0sb} +
C^c_{0sb}\cos\phi_k + C^s_{0sb}\sin\phi_k]$, and for Z100 the rotated ACS
coefficients are used.  The marginal over $\phi_i$ is then~\eqref{eq:lse}.

\subsection{Full Log-Likelihood}

Since shots are independent given $\beta$:
\begin{equation}
  \log\Lcal(\beta)
  = \sum_{i=1}^{N_\mathrm{shots}} \log\Lcal_i^\mathrm{marg}(\beta).
  \label{eq:full_ll}
\end{equation}

%%============================================================
\section{Implementation}
\label{sec:impl}
%%============================================================

\subsection{ACS Pre-Computation (Once per Shot per AI)}

The following steps are performed once per shot per AI and the results cached
to disk as a compressed NumPy archive (\texttt{.npz}) keyed by
\texttt{run/site/bins/n\_quad/K}:

\begin{enumerate}

  \item \textbf{Cloud nuisance profiling.}
        Solve~\eqref{eq:eta_profile} with L-BFGS-B (up to 200 iterations,
        \texttt{ftol}$=10^{-10}$, \texttt{gtol}$=10^{-8}$, 20 line-search steps).
        The objective evaluates the logsumexp over $K$ phase-grid points for
        both ports jointly at the current $\eta$ via \texttt{SurrogatePixelACS}.

  \item \textbf{ACS at $\hat\eta_{iz}$.}
        Call \texttt{SurrogatePixelACS.pixel\_acs}$(\hat\eta_{iz})$ to obtain
        tensors $\{A_{zsb}, C^c_{zsb}, C^s_{zsb}\}_{s \in \{g,e\}}$ of shape
        $(N_b,)$ each, on GPU.

  \item \textbf{Scale factor.}
        Compute $L_{iz}$ via~\eqref{eq:Lfix}.

  \item \textbf{Cache.}
        Store ACS tensors, $L_{iz}$, $\hat\eta_{iz}$, and raw pixel counts
        $(n_{izgsb}, n_{izesb})$ to \texttt{.npz}.

\end{enumerate}

\subsection{Signal Fit (Per Run)}

Given the cached ACS tables for all shots in a run, fit $\beta = (A_s, A_c)$:

\begin{enumerate}

  \item \textbf{Signal phases.}
        $\delta\phi_i = A_s\sin(2\pi f\,t_i) + A_c\cos(2\pi f\,t_i)$
        for $i = 1,\ldots,N_\mathrm{shots}$.

  \item \textbf{Rotation trick.}
        For each shot $i$, rotate the Z100 ACS tensors by $\delta\phi_i$.

  \item \textbf{Expected counts.}
        For each phase point $\phi_k$, compute $\lambda_{izsb}(\phi_k)$ via
        \eqref{eq:lambda_analytic} as a GPU broadcast:
        \texttt{lam[k, b] = L * max(A[b] + Cc[b]*cos\_phi[k]
        + Cs[b]*sin\_phi[k], 1e-300)}.

  \item \textbf{Conditional logL.}
        Compute~\eqref{eq:ll_phik} for each $(i, k)$ pair.

  \item \textbf{Phase marginalisation.}
        Apply logsumexp over $k$ and subtract $\log K$.

  \item \textbf{Shot sum and optimisation.}
        Sum over shots; negate and pass to L-BFGS-B with bounds
        $A_s \in [0, \pi]$, $A_c \in [-\pi, \pi]$.

\end{enumerate}

\subsection{Multi-Start Strategy for $\beta$}

The marginal log-likelihood $\log\Lcal(A_s, A_c)$ is smooth but can be flat
near the null $(0,0)$ when the signal is weak.  We use multi-start L-BFGS-B
from amplitude starting values $A_s^{(0)} \in \{0, 0.03, 0.07, 0.10, 0.15\}$
at $A_c^{(0)} = 0$, take the best solution, and report convergence diagnostics.

\subsection{Computational Cost}

\begin{table}[h]
\centering
\begin{tabular}{lll}
\toprule
Step & Phase & Typical wall-time \\
\midrule
\texttt{SurrogatePixelACS} construction   & Once (shared) & $\sim\!12$\,s \\
Cloud profiling ($K{=}512$ logsumexp)     & Per shot/AI   & $\sim\!0.5$--$2$\,s \\
All shots profiling + ACS cache           & Per run       & $\sim\!1$--$2$\,min \\
Signal fit (50 shots, multi-start)        & Per run       & $\sim\!40$\,s \\
\bottomrule
\end{tabular}
\caption{Typical timings on one GPU (A100/V100 class) with $K=512$, $N_b=1024$
($32\times32$ bins), $n_\mathrm{quad}=20{,}000$, $N_\mathrm{shots}=50$.
Cloud profiling timings dominate on the first pass; subsequent runs load from
the \texttt{.npz} cache in milliseconds.}
\label{tab:cost}
\end{table}

%%============================================================
\section{Maximum Likelihood Estimation}
%%============================================================

\subsection{Science Parameters}

With $f$ fixed, the only free parameters are $\beta = (A_s, A_c)$.
The MLE is
\begin{equation}
  \hat\beta = \arg\max_\beta\; \log\Lcal(\beta),
  \label{eq:mle}
\end{equation}
solved by multi-start L-BFGS-B (Section~\ref{sec:impl}).  The recovered
amplitude and phase are $\hat A = \sqrt{\hat A_s^2 + \hat A_c^2}$ and
$\hat\varphi = \arctan2(\hat A_c, \hat A_s)$.

\subsection{True Signal Parameterisation}

The simulated signal is generated as
$\delta\phi_i = A_\mathrm{true}\sin(2\pi f t_i + \varphi_\mathrm{true})$.
Expanding via angle addition:
\begin{equation}
  A_s^\mathrm{true} = A_\mathrm{true}\cos\varphi_\mathrm{true},
  \qquad
  A_c^\mathrm{true} = A_\mathrm{true}\sin\varphi_\mathrm{true},
\end{equation}
so the ground-truth $(A_s, A_c)$ is read directly from the H5 file attributes
\texttt{signal\_amp} and \texttt{signal\_phase}.  Note that $\delta\phi_i$ is
stored only in the Z100 H5 file (it is zero in Z0 by construction).

%%============================================================
\section{Full Per-Bin IS Marginalisation}
\label{sec:full_bin_is}
%%============================================================

This section describes the pipeline implemented in
\texttt{python-scripts/full\_image\_is\_inference.py}.  It uses the
complete spatial image (all $N_b$ per-bin Poisson counts) without
profiling the cloud nuisances per shot.  Cloud parameters $\eta$ are
marginalised by importance sampling (IS) from the finite-$R$ moment
posterior $p(\eta \mid m_i)$.

\subsection{Estimator}

The marginal shot likelihood is
\begin{equation}
  \mathcal{L}_i(\beta)
  = \int p(\text{image}_i \mid \beta, \eta)\; p(\eta)\; d\eta
  \;\approx\;
  \frac{1}{M}\sum_{j=1}^M w_j\; p(\text{image}_i \mid \beta, \eta_j),
  \label{eq:is_full_estimator}
\end{equation}
where $\eta_j \sim q(\eta)$ is drawn from the IS proposal and
\begin{equation}
  w_j = \frac{p(\eta_j)}{q(\eta_j)}
  \label{eq:is_weight_full}
\end{equation}
corrects for the mismatch between $q$ and the prior $p(\eta)$.  The
weights are functions of $\eta$ only — they carry no $\beta$-dependence
— so the same samples and weights are reused across all $\beta$
evaluations.

\subsection{IS Proposal: $\alpha$-Inflated Moment Posterior}

The proposal is the moment-conditioned posterior with covariance inflated
by $\alpha > 1$:
\begin{equation}
  q(\eta) = p(\eta \mid m_i,\alpha) = \mathcal{N}(\hat\eta_i,\, \alpha\, P_{\rm post}),
  \label{eq:proposal}
\end{equation}
where $\hat\eta_i$ and $P_{\rm post}$ come from a finite-$R$ Kalman
update on the four image moment observations
$(\hat\mu_{xf},\hat\mu_{yf},\hat\sigma^2_{xf},\hat\sigma^2_{yf})$.
Unlike the R=0 update, the finite-$R$ posterior is a non-degenerate 8D
Gaussian, so sampling from it is straightforward.

The 8D posterior factorises into four independent 2D blocks:
\begin{itemize}
  \item \textbf{Mean-$x$ block} $(\mu_{x_0}, \mu_{v_x}) \mid \hat\mu_{xf}$:
    observation $\hat\mu_{xf}$ with noise variance $R_{\mu,x} = \hat\sigma^2_{xf}/N_i$,
    observation matrix $h = (1,\, T)$.
  \item \textbf{Mean-$y$ block}: identical with $y$ moments.
  \item \textbf{Width-$x$ block} $(\sigma_{x_0}, \sigma_{v_x}) \mid \hat\sigma^2_{xf}$:
    linearised observation $h_\sigma = (2\bar\sigma_{\rm pos},\, 2T^2\bar\sigma_{\rm vel})$
    (delta-method gradient of $\hat\sigma^2_{xf}$ w.r.t.\ $\sigma_{x_0}, \sigma_{v_x}$),
    noise variance $R_{V,x} = 2\hat\sigma^4_{xf}/N_i$.
  \item \textbf{Width-$y$ block}: identical with $y$ moments.
\end{itemize}
Each block is a standard 2D Kalman update with scalar observation.  The
four posterior covariance matrices are assembled into a block-diagonal
$P_{\rm post}\in\mathbb{R}^{8\times8}$.

Inflation by $\alpha > 1$ (default $\alpha = 4$) ensures $q$ has heavier
tails than the true posterior, so importance weights $w_j$ remain bounded
with high probability.

\subsection{Per-Bin ACS via 2D Gauss-Hermite}

For each IS sample $\eta_j$, the per-bin ACS tensors
$(A_{zsb}^{(j)}, C^c_{zsb}{}^{(j)}, C^s_{zsb}{}^{(j)})$ of shape
$(N_b,)$ are computed via \texttt{SemiAnalyticPixelACS.pixel\_acs}.
This uses:
\begin{itemize}
  \item \emph{Exact spatial weight} $P_b = \Phi(x_{b+}/\sigma_{xf}) -
        \Phi(x_{b-}/\sigma_{xf})$ via Gaussian CDF per bin.
  \item \emph{2D GH quadrature} over the conditional velocity
        distribution $p(v_{x_0},v_{y_0} \mid x_f=x_b, y_f=y_b)$.
        For each bin~$b$, the conditional velocity mean and variance
        follow from standard Gaussian conditioning:
        \begin{equation}
          v_{x_0} \mid x_b
          \;\sim\;
          \mathcal{N}\!\Bigl(
            \mu_{v_x} + \frac{T\sigma^2_{v_x}}{\sigma^2_{xf}}(x_b - \mu_{xf}),\;
            \frac{\sigma^2_{v_x}\sigma^2_{x_0}}{\sigma^2_{xf}}
          \Bigr),
        \end{equation}
        giving $n_{\rm gh}^2$ deterministic quadrature nodes per bin.
\end{itemize}
For $M=64$ IS samples, $N_b=400$ bins and $n_{\rm gh}=4$ (16 nodes per
bin), the precompute requires 64 sequential calls to
\texttt{pixel\_acs}, each involving $400\times16 = 6400$ GPU
evaluations of the PSMAP surrogate.

\subsection{Per-Bin Poisson Log-Likelihood}

The data model for a single AI uses all $N_b$ bins:
\begin{equation}
  p(\text{image}_i \mid \beta, \eta_j, \phi)
  = \prod_{b,s} \mathrm{Poisson}(n_{izsb};\;\lambda_{izsb}^{(j)}(\phi)),
  \label{eq:poisson_per_bin}
\end{equation}
where $\lambda_{izsb}^{(j)} = \Lambda_{iz}^{(j)}(\phi)
[A_{zsb}^{(j)} + C^c{}_{zsb}^{(j)}\cos\phi + C^s{}_{zsb}^{(j)}\sin\phi]$
and the normalisation $\Lambda_{iz}^{(j)}(\phi)$ is chosen so that the
total predicted count equals the observed total $N_{iz}$ (as in
\eqref{eq:Lphi}).

The per-bin Poisson log-likelihood for IS sample $j$ at phase $\phi_k$ is
\begin{equation}
  \ell_{iz}^{(j)}(\phi_k)
  = \sum_b \Bigl[
      n_{izgb}\log\lambda_{izgb}^{(j)}(\phi_k)
      - \lambda_{izgb}^{(j)}(\phi_k)
      + n_{izeb}\log\lambda_{izeb}^{(j)}(\phi_k)
      - \lambda_{izeb}^{(j)}(\phi_k)
    \Bigr].
  \label{eq:ll_per_bin_is}
\end{equation}
Summing over bins is accumulated in chunks of size \texttt{chunk\_bins}
to bound peak memory to $\mathcal{O}(M \times \texttt{chunk\_bins} \times K)$
instead of $\mathcal{O}(M \times N_b \times K)$.

\subsection{IS-Weighted Phase Marginalisation}

The IS-marginalised log-likelihood for AI $z$ at phase $\phi_k$ is
\begin{equation}
  \log F_{iz}(\phi_k)
  = \operatorname{logsumexp}_{j=1}^M
    \Bigl[\log w_j + \ell_{iz}^{(j)}(\phi_k)\Bigr]
    - \operatorname{logsumexp}_{j=1}^M \log w_j,
  \label{eq:logF_iz_is}
\end{equation}
which normalises the IS weights.  The phase marginalisation and shot
combination proceed as in \eqref{eq:logLi} and \eqref{eq:logLfull_gh}.

\subsection{ESS Diagnostic}

The effective sample size
\begin{equation}
  \mathrm{ESS}
  = \frac{(\sum_j w_j)^2}{\sum_j w_j^2}
  \;\in\; [1, M]
  \label{eq:ess_full}
\end{equation}
is reported per shot per AI.  $\mathrm{ESS}/M \approx 1$ indicates that
the full spatial image adds no cloud information beyond the moments
(the moment posterior $q$ is already the correct posterior).
$\mathrm{ESS}/M \ll 1$ indicates that the PSMAP fringe structure
constrains $\eta$ beyond the cloud envelope moments.

\subsection{Default Hyperparameters}

\begin{table}[h]
\centering
\begin{tabular}{lll}
\toprule
Parameter & Symbol & Default \\
\midrule
Inference bins per side          & $\sqrt{N_b}$ & 20 \\
GH order per velocity dimension  & $n_{\rm gh}$ & 4 \\
IS samples per shot per AI       & $M$ & 64 \\
Proposal inflation factor        & $\alpha$ & 4.0 \\
Phase-grid points                & $K$ & 128 \\
Chunk size (logL memory)         & \texttt{chunk\_bins} & 64 \\
\bottomrule
\end{tabular}
\caption{Default hyperparameters for \texttt{full\_image\_is\_inference.py}.
With these defaults, peak RAM for the logL computation is
$M \times \texttt{chunk\_bins} \times K \times 3 \times 8 \approx 6$\,MB.}
\label{tab:full_is_params}
\end{table}

%%============================================================
\section{Validation}
%%============================================================

Running \texttt{fit\_surrogate\_profiled\_signal.py} with oracle cloud
initialisation (\texttt{--init-from-true}) on the default 1\,M-atom dataset
($N_\mathrm{shots}=50$, $K=512$, $n_\mathrm{quad}=20{,}000$) gives:

\begin{itemize}
  \item Signal amplitude RMSE $\approx 0.0013$\,rad
        (1.3\,\% of true $A=0.1$\,rad).
  \item Signal phase RMSE $\approx 0.009$\,rad.
  \item Cloud COM RMSE: $\sim\!2$--$5$\,$\mu$m for position,
        $\sim\!1$--$4$\,$\mu$m/s for velocity.
  \item Cloud width RMSE: $\sim\!0.1$--$0.3$\,$\mu$m.
  \item L-BFGS-B converges in $\sim\!5$ iterations and $\sim\!25$ function
        evaluations for both the cloud profiling and the signal fit.
\end{itemize}

%%============================================================
\section{\texttt{map\_inference.py}: $\phi$ Multistart and Empirical RMSE
vs.\ CRB (2026-07-16)}
\label{sec:map_inference_rmse}
%%============================================================

This section documents a direct empirical test of \texttt{map\_inference.py}
(the plug-in-MAP, moment-based cloud estimator\footnote{See
the module docstring: $\eta_i$ is a closed-form constrained-MAP estimate from
observed final-image moments (Kalman update, $R=0$); this is \emph{not} an
iterative nonlinear optimisation, so there is no local-minima risk for
$\eta$ itself.}): does hardening $\phi$'s per-shot solve with genuine
multistart change $\hat\beta$'s RMSE, and how does the achieved RMSE compare
to the Cram\'er-Rao bound, across two atom counts and multiple $\eta$-handling
strategies?

\subsection{$\phi$ Multistart: No Measurable Effect}

\texttt{solve\_phi\_batch} originally seeded its Newton polish from only the
single best point of a coarse scan, relying on an unverified assumption
(stated in its own docstring) that the port-summed combined likelihood is
``effectively unimodal.'' This was hardened to genuine multistart: locate the
top-$K$ local maxima of the coarse scan per shot, Newton-polish \emph{each},
keep whichever converges highest (implemented via a new
\texttt{phi\_n\_multistart} parameter; $\ge 2$ triggers the new behaviour,
$=1$ recovers the original code path exactly).

Tested across all 80 real independent runs of the $10^6$-atom dataset
(\texttt{moments} mode, $N=200$ shots, \texttt{gh\_order=12},
\texttt{phi\_method=point}): $\hat\beta$'s RMSE was \textbf{identical to 5
decimal places} between \texttt{phi\_n\_multistart=1} and \texttt{=6}
($\mathrm{RMSE}(A_s)=0.00019$, $\mathrm{RMSE}(A_c)=0.00017$ both ways;
per-run $\hat\beta$ differences $\le 2\times10^{-6}$, nonzero for only 14/80
runs). \textbf{Conclusion: for this port-summed statistic, the original
single-scan-point assumption holds up empirically, unlike the
pixel-resolved case} (\texttt{joint\_profile\_inference.py}, where $\phi$ is
100\% multimodal across a 50-shot sample -- see
\texttt{notes/joint\_profile\_inference.tex}, Issue 4). Multistart is kept
available (cheap relative to the rest of the pipeline) but is not required
for correctness here.

\subsection{Three-Way $\eta$-Handling Comparison, $10^6$ Atoms, $N=200$ Shots}

80 real independent runs, \texttt{phi\_method=point} (pure point-estimate
profiling, no Laplace correction -- this isolates profiling's own accuracy
from any marginalisation approximation), \texttt{quad=gh},
\texttt{gh\_order=12} (chunked over shots, \texttt{gh\_chunk\_shots=30}, to
fit GPU memory -- see \texttt{--gh\_chunk\_shots} in the CLI; unchunked
$n\_order{=}12$ at 200 shots OOMs a 46\,GB GPU, since the tensor-product GH
rule here is \emph{4D} -- over the cloud's full initial
$(x_0,y_0,v_{x0},v_{y0})$ state, integrated unconditionally, as opposed to
the \emph{2D} conditional-velocity integral (\texttt{pixel\_n\_gh} in
\texttt{joint\_profile\_inference.py}) used for the per-pixel conditional
likelihood elsewhere in this codebase -- the two are not directly comparable
quadrature orders):

\begin{table}[h]
\centering
\begin{tabular}{lccc}
\toprule
Mode & RMSE($A_s$) & RMSE($A_c$) & vs.\ moments \\
\midrule
\texttt{use\_moments=0} (prior-mean $\eta$, no position info) & 0.00051 & 0.00049 & 2.7$\times$ worse \\
\texttt{use\_moments=1} (moment-based $\eta$, realistic) & 0.00019 & 0.00017 & --- \\
\texttt{use\_true\_eta=1} (oracle $\eta$) & 0.00015 & 0.00015 & 20--25\% better \\
\bottomrule
\end{tabular}
\caption{$\hat\beta$ RMSE across 80 independent $10^6$-atom runs, $N=200$
shots. All three modes are essentially unbiased (bias $\le 8\times10^{-5}$).}
\end{table}

Interpretation: moment-based $\eta$ estimation captures \emph{most} of the
available gain from spatial information -- the jump from no-info to moments
is large (2.7$\times$), but moments to perfect oracle knowledge only buys a
further $\sim$20--25\%.

\textbf{A CRB-reference caveat.} The natural first comparison is against the
port-summed CRB (\texttt{crb\_signal.py}, \texttt{likelihood=port}:
$\sigma=(0.00022, 0.00021)$ at $N=200$) -- \texttt{use\_moments=1} beats
this. But this is a \emph{mismatched} reference: the port CRB's Fisher
information comes from the port-summed Poisson statistic's own dependence on
$\eta$ (properly Schur-complementing $\eta$'s prior uncertainty), not from
explicit position moments. \texttt{use\_moments=1} uses a strictly richer
information channel (explicit $(\mu_{xf},\mu_{yf},V_{xf},V_{yf})$ moments)
than the port CRB assumes, so beating it is expected, not a notable
achievement in itself. The three-way comparison above (moments vs.\ no-info
vs.\ oracle) is the more meaningful, apples-to-apples test.

\subsection{Matched $10^6$ vs.\ $10^8$ Atom Comparison, $N=50$ Shots}

The $10^8$-atom dataset only has runs available at $N=50$ shots/run (40
independent runs), so a fair atom-count comparison requires truncating the
$10^6$-atom dataset to the same $N=50$ (\texttt{--max\_shots 50}) rather than
comparing its $N{=}200$ results directly -- CRB/RMSE scale with
$1/\sqrt{N_\mathrm{shots}}$ independent of atom count, so mismatched $N$
would confound the comparison.

\begin{table}[h]
\centering
\begin{tabular}{lcccc}
\toprule
Config & RMSE($A_s$) & RMSE($A_c$) & Port CRB & Oracle-$\theta$ CRB \\
\midrule
$10^6$, moments & 0.00049 & 0.00045 & (0.000434, 0.000477) & (0.000159, 0.000157) \\
$10^6$, oracle-$\eta$ & 0.00034 & 0.00035 & --- & (0.000159, 0.000157) \\
$10^8$, moments & 0.00029 & 0.00024 & (0.000287, 0.000425) & (0.0000160, 0.0000160) \\
$10^8$, oracle-$\eta$ & 0.00003 & 0.00003 & --- & (0.0000160, 0.0000160) \\
\bottomrule
\end{tabular}
\caption{Matched-$N{=}50$ comparison, port-summed statistic. Oracle-$\theta$
CRB uses the theta-known Fisher block (\texttt{H\_pp\_no\_prior\_no\_theta}
in \texttt{crb\_signal.py}'s per-shot diagnostics), the correct matched
bound for \texttt{use\_true\_eta=1} (theta uncertainty eliminated, not just
estimated well). Pixel-likelihood CRB at the same $N=50$, for reference:
$10^6$: $\sigma=(0.000162,0.000160)$, oracle$=(0.000159,0.000156)$;
$10^8$: $\sigma=(0.000018,0.000018)$, oracle$=(0.0000160,0.0000160)$ --
pixel CRB sits very close to its own oracle at both atom counts (the full
image already constrains $\theta$ almost as well as knowing it exactly),
unlike the port statistic, whose CRB is far from its oracle -- consistent
with the $\sim\!96\%$ vs.\ $\sim\!9\%$ phase-information comparison
motivating this whole line of work (\S1).}
\end{table}

Two findings worth flagging plainly, neither yet explained:

\begin{itemize}
\item \textbf{The oracle-vs-oracle-CRB gap is consistent across a
$100\times$ change in photon count}: $10^6$ oracle-$\eta$ RMSE is
$2.2\times$ its matched CRB; $10^8$ oracle-$\eta$ RMSE is $1.9\times$ its
matched CRB. This consistency across regimes with very different Poisson
noise scales argues against a simple photon-counting or discretisation
explanation and points toward something structural in the estimation
pipeline -- \texttt{solve\_phi\_batch}'s convergence precision
(\texttt{n\_scan}, \texttt{n\_newton}, finite-difference step $h$) and
whether point-estimate profiling is unbiased-but-inefficient (as opposed to
achieving the CRB) are both unverified candidate explanations, not
established causes. \textbf{Not yet root-caused.}
\item \textbf{Atom-count scaling is far weaker than naive photon-counting
would predict}: $10^6\to10^8$ atoms (100$\times$ more photons/shot) only
improves moments-mode RMSE by $\sim\!1.7\times$, not
$\sqrt{100}=10\times$. This suggests the $\eta$ prior's own width -- not
per-shot photon-counting precision -- is the effective floor for the
moments-based estimator at these atom counts. (The oracle-$\eta$ estimator,
by contrast, does show the expected $\sim\!10\times$ improvement:
$0.00034\to0.00003$, $0.00035\to0.00003$ -- consistent with this being a
property of the moment-based $\eta$ \emph{estimation} step specifically,
not the underlying photon statistics.)
\end{itemize}

\end{document}
